\begin{titlepage}
  \sffamily
  \vspace*{30mm}
  {\color{brand}\rule{\linewidth}{1.5pt}}\\[8mm]
  {\Huge\bfseries\color{brand} Aconcagua CFD}\\[4mm]
  {\LARGE User Guide}\\[6mm]
  {\large Drag of the Aconcagua rocket with OpenFOAM v2412:\\
   geometry, mesh, cases, sweeps, measurement and validation}\\[8mm]
  {\color{brand}\rule{\linewidth}{1.5pt}}
  \vfill
  \begin{tabular}{@{}ll@{}}
    Document & ACG-CFD-UG-001 \\
    Revision & \docrev \\
    Date & \docdate \\
    Repository & \texttt{github.com/vnemegabriel/CFD-ITBA-ROCKETRY} \\
    Applies to & branch \texttt{sweep/nemepad}, \texttt{flightModel/} \\
    Organization & ITBA Rocketry \\
  \end{tabular}
  \vspace*{10mm}
\end{titlepage}

\pagenumbering{roman}
\chapter*{Document control}
\addcontentsline{toc}{chapter}{Document control}
\markboth{Document control}{Document control}

\section*{Revision history}

\begin{tabularx}{\linewidth}{@{}l l L@{}}
\toprule
Rev. & Date & Change \\
\midrule
A & \docdate & First issue. Replaces the Markdown files \file{WORKFLOW.md}, \file{MESH.md},
\file{SOLVERS.md}, \file{VALIDATION.md} and \file{MEASUREMENTS.md} of \file{flightModel/docs/}.
Adds the chapters on the comparison with Nektar++, the flight data and troubleshooting. \\
\bottomrule
\end{tabularx}

\section*{Status of the template}

\begin{caution}
The case templates are not validated. Do not publish a result before you do the
procedures in \cref{ch:validation}. The supersonic template has an open defect
(\cref{sec:ts-central-divergence}).
\end{caution}

\section*{Maintenance of this document}

\begin{itemize}
\item This guide is the only user documentation of \file{flightModel/}. Do not
  write the same information in another file.
\item Change the source in \file{flightModel/docs/guide/} and compile it with
  \cmd{make} (\cref{sec:build-guide}). Commit the source and the PDF together.
\item When a change to a script changes a procedure, change the procedure in this
  guide in the same commit.
\item Increase the revision letter and add a line to the revision history for each
  issue.
\end{itemize}

\tableofcontents
\listoffigures
\listoftables

\chapter*{About this guide}
\addcontentsline{toc}{chapter}{About this guide}
\markboth{About this guide}{About this guide}

\section*{Purpose}

This guide tells you how to calculate the drag coefficient of the Aconcagua rocket
as a function of the Mach number with OpenFOAM. It gives the procedures, the
reference data of each dictionary, the models, the measurements that each run
records, and the procedures that validate the results.

\section*{Audience}

The guide is for members of ITBA Rocketry who run or change the CFD cases. You must
know how to use a Linux shell. You must know the basic concepts of compressible
flow, RANS turbulence models and the finite-volume method. You do not have to know
the OpenFOAM source code.

\section*{Scope}

The guide covers the directory \file{flightModel/} of the repository: the geometry
script, the mesh scripts, the two case templates, the sweep scripts, the
post-processing scripts and the flight data. It does not cover the directory
\file{fin_flutter/}.

\section*{How to use this guide}

\begin{tabularx}{\linewidth}{@{}P{0.38\linewidth} L@{}}
\toprule
If you want to & Read \\
\midrule
Make the first run & \cref{ch:overview,ch:setup}, then \cref{sec:make-case,sec:run-case} \\
Run the sweep of record & \cref{ch:sweeps} \\
Change the geometry or the mesh & \cref{ch:geometry,ch:mesh} \\
Know why the solver is set up as it is & \cref{ch:solvers} \\
Find a result or a measurement & \cref{ch:results} \\
Compare a run with Nektar++ & \cref{ch:nektar} \\
Compare a run with the flight & \cref{ch:flight} \\
Validate a result & \cref{ch:validation} \\
Correct a problem & \cref{ch:troubleshooting} \\
Find a command quickly & \cref{app:quickref} \\
\bottomrule
\end{tabularx}

\section*{Conventions}

\begin{tabularx}{\linewidth}{@{}P{0.34\linewidth} L@{}}
\toprule
Convention & Meaning \\
\midrule
\file{mesh/Allmesh} & A file, a directory, a command, a dictionary entry or a value. \\
\file{<mesh-dir>} & A value that you must supply. \\
\file{[stl-file]} & An optional argument. \\
\file{~/meshes}, \file{~/runs} & The default directories of the meshes and of the runs. \\
Paths without a leading \file{/} & Relative to \file{flightModel/}, except where the text says otherwise. \\
$h0$, $h1$, $h2$, $h3$ & The four meshes of the sweep (\cref{tab:meshes}). \\
\texttt{m08-a04-h1} & A run: Mach 0.8, angle of attack $4^\circ$, mesh $h1$. \\
\bottomrule
\end{tabularx}

\medskip
Procedures are numbered steps. Each step gives one action. A result that you must
see after a step follows the step.

\section*{Notices}

\begin{caution}
A CAUTION tells you about a condition that can damage data, give an incorrect
result without a message, or waste days of computer time.
\end{caution}

\begin{note}
A NOTE gives information that helps you to do a step or to understand a result.
\end{note}

\section*{Writing standard}

The structure of this guide follows ISO/IEC/IEEE 26514 \citep{iso26514}: concepts,
procedures and reference information are in different sections. The text follows
the writing rules of ASD-STE100 Simplified Technical English \citep{asdste100}:
short sentences, one instruction in each sentence, the active voice and one meaning
for each word. Technical names (\emph{mesh}, \emph{solver}, \yplus,
\emph{wall function}) stay as the OpenFOAM documentation writes them.

\section*{Related documents}

\begin{tabularx}{\linewidth}{@{}P{0.38\linewidth} L@{}}
\toprule
Document & Contents \\
\midrule
\file{CLAUDE.md} & Rules for the people and the agents who change the repository. \\
\file{flightModel/case-central/README.md}, \file{case-pimple/README.md} & Short description of each template, with references to this guide. \\
\file{flightModel/flightData/cdFlight.py} & Source of the flight analysis of \cref{ch:flight}. \\
OpenFOAM v2412 user guide & \url{https://www.openfoam.com/documentation/user-guide} \\
\bottomrule
\end{tabularx}
